\section*{Chapter \ref{ch:b3:growth-empirics} --- Growth Empirics and Convergence}

\begin{solution}{exo:b3:growth-empirics:1}
The 1980s, with a median of 1.2~\% a year; 37~\% of countries saw income per person fall. % figdata: decades_stats
\end{solution}

\begin{solution}{exo:b3:growth-empirics:2}
$\lambda = -\ln(1 - 0.015 \times 40)/40 = -\ln 0.4/40 = 2.3~\%$ a year.
\end{solution}

\begin{solution}{exo:b3:growth-empirics:3}
$\ln 2/0.02 = 34.7$ years and $\ln 2/0.07 = 9.9$ years.
\end{solution}

\begin{solution}{exo:b3:growth-empirics:4}
Poor countries with the same saving, schooling and population growth as rich ones
growing no faster, or more slowly, once those are held fixed --- a zero or positive
slope on initial income in the conditional regression.
\end{solution}

\begin{solution}{exo:b3:growth-empirics:5}
The tested solution is \code{ex05_convergence.py}, in the chapter's \code{exercises/}
folder. All 111 countries: slope $-0.0011$; the 25 richest in 1960: $-0.0015$, hardly any
convergence; the 25 richest in 2019: $-0.0159$, ten times as much. The last sample is chosen on the outcome: countries that were poor in 1960
enter it only if they grew fast, and those that were rich and fell behind drop out,
so convergence appears by construction.
\end{solution}

\begin{solution}{exo:b3:growth-empirics:6}
The tested solution is \code{ex06_persistence.py}, in the chapter's \code{exercises/}
folder. Growth correlations between successive decades are $+0.45$, $+0.44$, $+0.33$, $-0.08$ and $+0.32$;
log levels correlate at 0.95 to 0.98. Levels persist, growth rates do not.
\end{solution}

\begin{solution}{exo:b3:growth-empirics:7}
The exponent is $1/(1 - 2/3) = 3$, so income is proportional to $s_k s_h/(n + g + \delta)^2$: the ratio is
$(0.24/0.12)(0.12/0.06) = 4$.
\end{solution}

\begin{solution}{exo:b3:growth-empirics:8}
Mean growth after is 5~\%, at least 3.5, and 4 points above the past mean of 1~\%; income keeps
rising, so the end year is a new peak: an acceleration is dated at the last year
of slow growth.
\end{solution}

\begin{solution}{exo:b3:growth-empirics:9}
A negative slope means the initially poor grew faster on average, but shocks of
large variance can spread incomes out at the same time
(\cref{prop:b3:growth-empirics:galton}); and a sample chosen after the fact, such as
today's rich countries, shows convergence by selection.
\end{solution}

\begin{solution}{exo:b3:growth-empirics:10}
With $\dot k = s k^{\alpha} - (n + g + \delta)k$, $d\ln k/dt = s e^{(\alpha - 1)\ln k} - (n + g + \delta)$. Its derivative with respect to $\ln k$ at $k^*$ is
$(\alpha - 1)s k^{*\alpha - 1} = -(1 - \alpha)(n + g + \delta)$, using $s k^{*\alpha - 1} = n + g + \delta$. So $d\ln k/dt \approx -(1 - \alpha)(n + g + \delta)(\ln k - \ln k^*)$, and since $\ln y = \alpha\ln k$, $\ln y$
converges at the same speed.
\end{solution}

\begin{solution}{exo:b3:growth-empirics:11}
The tested solution is \code{ex11_bounds.py}, in the chapter's \code{exercises/}
folder. Investment has bounds $[+0.54, +1.53]$ and is robust; every other candidate is fragile,
including schooling, whose true effect is 0.3 but whose bounds are $[-0.05, +0.78]$: its correlation
with investment makes its estimate imprecise when investment enters.
\end{solution}

\begin{solution}{exo:b3:growth-empirics:12}
$0.4 \times 1 + 0.6 \times 3 = 2.2$ becomes $0.3 \times 1 + 0.7 \times 3 = 2.4$: a rise of 9.1~\%. With 1.5 outside farming, 1.3 becomes 1.35: 3.8~\%.
\end{solution}

\begin{solution}{pb:b3:growth-empirics:fifty}
% figdata: brenland_stats
\textbf{1.} 12{,}377, 24{,}675 and 31{,}470 dalens of 2015.

\textbf{2.} 2.76~\% and 0.97~\% a year.

\textbf{3.} $\ln 2/0.00973 = 71$ years.

\textbf{4.} $31{,}470/12{,}377 = 2.54$.

\textbf{5.} Fourteen years at 4~\% multiply income by $e^{0.56} = 1.75$; growth compounds.

\textbf{6.} $\beta = -0.0115$.

\textbf{7.} $-\ln(1 - 50 \times 0.0115)/50 = 1.7~\%$ a year.

\textbf{8.} 2.04~\% a year; Brenland grew 1.87~\%, 0.17 points less.

\textbf{9.} Absolute in form, but the peers share one generating rule --- one steady
state --- so it is convergence within a group with common determinants.

\textbf{10.} 0.51 in 1975 and 0.28 in 2025. With a slope of $-0.0115$ over 50 years, $1 + 50\beta = 0.42$ shrinks the initial
spread, more than the noise of 0.4 points a year (0.2 in logs over 50 years) adds: here \omterm{def:b3:growth-empirics:convergence}{beta
convergence} comes with \omterm{def:b3:growth-empirics:convergence}{sigma convergence}.

\textbf{11.} Measurement error in 1975 income enters both sides and biases $\beta$ towards
convergence; so does omitting steady-state differences correlated with initial
income in the other direction.

\textbf{12.} 1980 and 1981 both qualify. For 1980 the next eight years average
$(0.5 + 7 \times 4)/8 = 3.56~\%$, a rise of 3.06; for 1981 the next eight average 4~\% against 0.5~\%
before, a rise of 3.5. The filter dates 1981, the last year of slow growth.

\textbf{13.} 0.5~\% a year in 1977--1981 and 4.0~\% in 1982--1989.

\textbf{14.} Yes: growth stays at 4~\% until 1995 and then 2~\%, still above the 0.5~\% before.

\textbf{15.} That the reform came before the take-off and was not a response to it;
that no other shock (\omterm{def:b3:ricardian-trade:terms-of-trade}{terms of trade}, a neighbour's boom) explains it; and how peers
that did not reform fared --- most reforms produce no acceleration.

\textbf{16.} The average over 1975--2025 mixes 0.5, 4, 2 and 1~\%; one average per country
cannot show when growth changed.

\textbf{17.} $0.2 \times 6 + 0.8 \times 16 = 14$ and $0.05 \times 12 + 0.95 \times 36 = 34.8$ thousand dalens.

\textbf{18.} $16/6 = 2.67$ and $36/12 = 3.0$.

\textbf{19.} Within $0.2 \times 6 + 0.8 \times 20 = 17.2$; between $-0.15 \times 6 + 0.15 \times 16 = 1.5$; cross $-0.15 \times 6 + 0.15 \times 20 = 2.1$; total 20.8.

\textbf{20.} The between term alone is 7.2~\% of the rise; with the cross term, 17.3~\%.

\textbf{21.} Part of the gap is measurement; and those who leave farming may be less
productive outside it than those already there.

\textbf{22.} Correcting for fewer hours and less schooling in farming lowers the gap,
and with it the gain from moving.

\textbf{23.} $0.2 \times 12 + 0.8 \times 36 = 31.2$, against 34.8: output per worker would have grown $\ln(31.2/14)/50 = 1.60~\%$ a year instead of
$\ln(34.8/14)/50 = 1.82~\%$.

\textbf{24.} Brenland converged with its peers at about 1.7~\% a year, and the move out of farming
accounts for about 17~\% of the rise in output per worker (7~\% at the initial sector productivities).
\end{solution}
